Large language models appear to generate rich, complex outputs from simple prompts, raising a fundamental question: does the information content of an LLM's output exceed that of its input?
We empirically investigate this question by measuring information content via two complementary compressors---gzip and LLM log-probabilities---across prompts of varying complexity, gibberish controls, and a paired human-vs-LLM text corpus.
We find that while LLM outputs contain more \emph{total} compressed bits than inputs (due to length expansion of $6$--$44\times$), the information \emph{density} per byte is consistently lower: outputs average $0.49\times$ the input density for complex prompts.
From the model's own perspective, outputs are nearly perfectly predictable at ${\sim}0.3$ \bpt, confirming that generation stays within the model's learned distribution.
We further show that the choice of compressor is decisive: gzip finds LLM text \emph{less} compressible than human text ($p < 10^{-11}$, Cohen's $d = -1.16$), while LLM-based compressors find the opposite---reconciling an apparent contradiction in prior work.
Gibberish inputs with high raw entropy produce outputs of normal information density, confirming that raw entropy does not equal usable information.
Our results support the view that LLMs act as \emph{decompressors}: they expand compressed semantic prompts into longer, predictable text, with the apparent information amplification sourced entirely from model weights.
